This appendix supplies the proofs that \S\ref{sec:lower_bound} defers. An
earlier pointer sent the reader to supplementary material; no such material
existed, and the authors correct that here.

The authors are explicit about what is fully proved and what is not.
Lemma~\ref{lem:id} and Lemma~\ref{lem:floor-general} are proved completely
below. Theorem~\ref{thm:general} is proved modulo one constant, and the authors isolate
exactly which step carries it rather than absorbing it silently.

\subsection{Proof of Lemma~\ref{lem:id}}

Fix any $w \in \R^d$ and recall $\bar H = \frac1T\sum_t H_t$ and that $\tauH$
is any solution of $\bar H \tauH = \frac1T\sum_t H_t \tau_t$ lying in
$\range(\bar H)$. Decompose $w - \tau_t = (w - \tauH) + (\tauH - \tau_t)$ and
expand the quadratic form:
\begin{align*}
  (w-\tau_t)^\top H_t (w-\tau_t)
  &= (w-\tauH)^\top H_t (w-\tauH) \\
  &\quad + 2 (w-\tauH)^\top H_t (\tauH - \tau_t)
   + (\tauH-\tau_t)^\top H_t (\tauH-\tau_t).
\end{align*}
Average over $t \in [T]$. The first term averages to
$(w-\tauH)^\top \bar H (w-\tauH)$ by linearity. The cross term averages to
\begin{equation*}
  2 (w-\tauH)^\top\Bigl[\bar H \tauH - \tfrac1T\sum_t H_t \tau_t\Bigr] = 0,
\end{equation*}
which vanishes by the defining property of $\tauH$ and is the only place that
property is used. The third term does not involve $w$. This gives
\eqref{eq:id-eq}. Since a maximum is at least an average,
$\max_t (w-\tau_t)^\top H_t (w-\tau_t)$ is at least the right-hand side.
\hfill$\square$

\paragraph{Two consequences used later.} First, any component of $w$
orthogonal to $\range(\bar H) = V_1 + \cdots + V_T$ annihilates against every
$H_t$, so a decoder may be assumed without loss to emit
$w \in \range(\bar H)$, a space of dimension $\deff$. Second, the bound
separates: a $w$-dependent compression term in the $\bar H$ metric, and a
$w$-independent floor.

\paragraph{An exact instance form.} Expanding the floor and using
$\frac1T\sum_t H_t\tau_t = \bar H\tauH$ and
$\tau_t^\top H_t \tau_t = B^2$ under $P^\star$,
\begin{equation}\label{eq:floor-exact}
  \frac1T\sum_t (\tau_t-\tauH)^\top H_t (\tau_t-\tauH)
  \;=\; B^2 - \tauH^\top \bar H \tauH .
\end{equation}
This form is worth stating because it is computable for any given cohort from
the adapters alone, with no appeal to $P^\star$. The authors use it in
\S\ref{sec:regimes}, and they note that it is the quantity a practitioner should
evaluate; the prior-averaged expression of Lemma~\ref{lem:floor-general} is a
statement about $P^\star$, not about the cohort in front of them.

\subsection{Proof of Lemma~\ref{lem:floor-general}}

The authors first record the second-moment identity that makes the result independent
of the spectra $D_t$.

\begin{claim}\label{clm:second-moment}
Under $P^\star$, for every $t$,
$\E[H_t \tau_t \tau_t^\top H_t \mid U_t] = (B^2/r)\, H_t$, and
$\tau_t^\top H_t \tau_t = B^2$ almost surely.
\end{claim}

\begin{proof}
With $H_t = U_t D_t U_t^\top$ and $\tau_t = U_t B D_t^{-1/2} u_t$, using
$U_t^\top U_t = I_r$,
\begin{equation*}
  H_t \tau_t = U_t D_t U_t^\top U_t B D_t^{-1/2} u_t
             = B\, U_t D_t^{1/2} u_t .
\end{equation*}
Hence $H_t \tau_t \tau_t^\top H_t = B^2 U_t D_t^{1/2} u_t u_t^\top D_t^{1/2}
U_t^\top$. Since $u_t \sim \mathrm{Unif}(S^{r-1})$ has
$\E[u_t u_t^\top] = I_r / r$, taking the conditional expectation gives
$(B^2/r)\,U_t D_t U_t^\top = (B^2/r) H_t$. For the second statement,
$\tau_t^\top H_t \tau_t = B^2 u_t^\top D_t^{-1/2} D_t D_t^{-1/2} u_t
= B^2 \|u_t\|^2 = B^2$. The $D_t^{-1/2}$ stretch in $P^\star$ is exactly what
makes both statements free of $D_t$.
\end{proof}

Now write $g := \frac1T\sum_t H_t \tau_t = \bar H \tauH$, so that by
\eqref{eq:floor-exact} the floor equals $B^2 - \tauH^\top \bar H \tauH$, and
since $\tauH = \bar H^+ g$ with $g \in \range(\bar H)$,
\begin{equation*}
  \tauH^\top \bar H \tauH
  = g^\top \bar H^+ \bar H \bar H^+ g = g^\top \bar H^+ g .
\end{equation*}
Condition on $\boldsymbol H$, equivalently on $(U_t, D_t)_{t}$. The $u_t$ are
independent across $t$ with $\E[u_t] = 0$, so $\E[H_t\tau_t] = 0$ and the
cross terms vanish:
\begin{equation*}
  \E[g g^\top \mid \boldsymbol H]
  = \frac{1}{T^2}\sum_{t=1}^T \E[H_t\tau_t\tau_t^\top H_t \mid \boldsymbol H]
  = \frac{1}{T^2}\sum_{t=1}^T \frac{B^2}{r} H_t
  = \frac{B^2}{rT}\,\bar H ,
\end{equation*}
using Claim~\ref{clm:second-moment}. Therefore
\begin{equation*}
  \E\bigl[g^\top \bar H^+ g \mid \boldsymbol H\bigr]
  = \tr\bigl(\bar H^+ \E[g g^\top \mid \boldsymbol H]\bigr)
  = \frac{B^2}{rT}\tr(\bar H^+ \bar H)
  = \frac{B^2}{rT}\rank(\bar H)
  = \frac{B^2 \deff}{Tr},
\end{equation*}
where $\tr(\bar H^+\bar H) = \rank(\bar H)$ because $\bar H^+\bar H$ is the
orthogonal projector onto $\range(\bar H)$, and
$\rank(\bar H) = \dim(V_1 + \cdots + V_T) = \deff$ since the $H_t$ are PSD, so
no cancellation can occur in the sum. Taking expectations over
$\boldsymbol H$ and substituting into \eqref{eq:floor-exact},
\begin{equation*}
  \E_{P^\star}\Bigl[\tfrac1T\sum_t (\tau_t-\tauH)^\top H_t(\tau_t-\tauH)\Bigr]
  = B^2 - \frac{B^2\,\E[\deff]}{Tr}
  = B^2\Bigl(1 - \frac{\E[\deff]}{Tr}\Bigr).
  \tag*{$\square$}
\end{equation*}

\paragraph{Remark: the floor is vacuous for LoRA cohorts in general
position.} The floor vanishes precisely when $\deff = Tr$, that is when the
$V_t$ are in direct sum. Under Stiefel-random sampling with $Tr \leq d$ this
holds almost surely, which is the regime the theorem's second display
assumes. It also holds for every real cohort the authors measured, including the
shared-initialisation ones whose subspaces are numerically near-collinear
(\S\ref{sec:regimes}). Near-collinear is not collinear: the $V_t$ remain
independent, $\deff = Tr$ exactly, and by \eqref{eq:floor-exact} the exact
floor is zero. What degrades in that regime is not the floor but the
conditioning of $\bar H$, and hence the norm of the interpolator achieving it.
The authors regard this as the main practical lesson of the lower bound and they develop
it in \S\ref{sec:regimes}.

\subsection{Proof of Theorem~\ref{thm:general}, and what it rests on}

By Yao's principle \eqref{eq:yao} it suffices to lower-bound the Bayes risk
under $P^\star$. By Lemma~\ref{lem:id}, for any code with output $w^\star$,
\begin{equation*}
  \E\bigl[\max_t D_t(w^\star)\bigr]
  \;\geq\;
  \underbrace{B^2\Bigl(1 - \frac{\E[\deff]}{Tr}\Bigr)}_{\text{Lemma
  \ref{lem:floor-general}}}
  \;+\;
  \E\bigl[(w^\star-\tauH)^\top \bar H (w^\star-\tauH)\bigr].
\end{equation*}
It remains to lower-bound the second term. As noted, the decoder may be
assumed to emit $w^\star \in \range(\bar H)$. Put
$\eta := \bar H^{1/2}\tauH$ and $\hat\eta := \bar H^{1/2} w^\star$, both in a
space of dimension $\deff$; the term equals $\E\|\hat\eta - \eta\|^2$, so the
subproblem is rate-$R$ quantisation of $\eta$ under squared error.

Condition on $\boldsymbol H$. Given $\boldsymbol H$ the only remaining
randomness is the directions $u_t$, and a conditional Shannon converse
therefore applies to every rate-$R$ code, \emph{including codebooks chosen
with knowledge of $\boldsymbol H$}. This is the step that makes the bound valid
against the $\boldsymbol H$-aware decoder of Appendix~\ref{app:rd-achv}: the authors do not
average over rotations, they condition. The standard converse for a
$\deff$-dimensional source under squared error gives
\begin{equation*}
  \E\bigl[\|\hat\eta - \eta\|^2 \bigm| \boldsymbol H\bigr]
  \;\geq\; \deff \cdot \frac{2^{2h(\eta \mid \boldsymbol H)/\deff}}
  {2\pi e}\cdot 2^{-2R/\deff},
\end{equation*}
with $h$ the differential entropy. From Lemma~\ref{lem:floor-general},
$\E[\|\eta\|^2 \mid \boldsymbol H] = B^2\deff/(Tr)$, so the per-coordinate
scale is pinned.

\paragraph{The constant, isolated.} This is the step, and the only step, that
uses Assumption~\ref{asm:epr}. Writing the entropy power of $\eta$ as a
constant fraction of its variance turns the display above into
\begin{equation*}
  c_{\mathrm{TQ}}\cdot \frac{B^2\,\E[\deff]}{Tr}\cdot 2^{-2R/\E[\deff]},
\end{equation*}
which is \eqref{eq:thm-lb}. The step the authors do not prove here is that the
entropy-power-to-variance ratio of $\eta$ is bounded below by an absolute
constant, uniformly in $(T, r, d)$ and in the spectra $D_t$. It is finite for
each fixed instance because $\eta$ has a density, and the scale is pinned by
$\tau_t^\top H_t\tau_t = B^2$; what the authors do not establish is uniformity. The authors
therefore quote $c_{\mathrm{TQ}}$ from TurboQuant
\citep{zandieh2025turboquant}, Theorem~3, which supplies exactly such a
constant for the rotation-invariant case, and they flag that transferring it to
their conditional setting is an assumption rather than a consequence of the
argument above.

\paragraph{Scope of the theorem.} With that caveat, Theorem~\ref{thm:general} should
be read as: the floor term is exact and unconditional (it is
Lemma~\ref{lem:floor-general}), and the rate term has the stated
\emph{form} $2^{-2R/\deff}$ with a constant the authors import rather than derive. The authors'
empirical work bears on the rate term only through
\S\ref{subsec:w3}, where the exponent is not resolved at the rates they could
sweep. Nothing in \S\ref{sec:regimes} depends on the constant.
